\documentclass{article}
\usepackage{amsmath,amssymb}
\usepackage[vocab={quadratics}]{leanify}
\newtheorem{lemma}{Lemma}
\begin{document}

\section{Notation}

We write $D(a,b,c)$ for the discriminant of the quadratic with coefficients $a$, $b$ and $c$.
We write $x_+(a,b,c)$ for the larger root of the quadratic with coefficients $a$, $b$ and $c$.
We write $x_-(a,b,c)$ for the smaller root of the quadratic with coefficients $a$, $b$ and $c$.
We write $V(a,b,c)$ for the vertex of the parabola with coefficients $a$, $b$ and $c$.
We write $N(a,b,c)$ for the number of real roots of the quadratic with coefficients $a$, $b$ and $c$.
We write $\Delta(p,q)$ for the discriminant of the cubic with coefficients $p$ and $q$.
We write $C(p,q)$ for the Cardano root of the cubic with coefficients $p$ and $q$.

\section{Quadratics}

\begin{lemma}\label{example}
$D(1,-3,2)=1$, $x_+(1,-3,2)=2$, $x_-(1,-3,2)=1$, $V(1,-3,2)=\frac32$ and $N(1,-3,2)=2$.
\end{lemma}

\begin{lemma}\label{vieta}
Let $a>0$, $b$ and $c$ be real numbers with $b^2\ge4ac$. Then $x_+(a,b,c)+x_-(a,b,c)=-\frac{b}{a}$.
\end{lemma}

\begin{lemma}\label{vieta-complex}
Let $a$, $b$ and $c$ be real numbers with $a\neq0$. Then the sum of the roots of the quadratic with coefficients $a$, $b$ and $c$ is $-\frac{b}{a}$, and the product of the roots of the quadratic with coefficients $a$, $b$ and $c$ is $\frac{c}{a}$.
\end{lemma}

\begin{lemma}\label{extreme}
The extreme value of the quadratic with coefficients $1$, $-2$ and $3$ is $2$, and the number of real roots of the quadratic with coefficients $1$, $0$ and $1$ is $0$.
\end{lemma}

\section{Cubics}

\begin{lemma}\label{cubic}
$\Delta(-3,2)=0$ and $C(0,-8)=2$.
\end{lemma}

\end{document}
